\documentclass{article}
\usepackage{iclr2025_conference,times}

\usepackage{graphicx}
\usepackage{hyperref}
\usepackage{url}

\iclrfinalcopy
\pagestyle{empty}

\title{Time-LLM: Time Series Forecasting by Reprogramming Large Language Models}
\author{\small\textbf{Ming Jin$^{1}$, Shiyu Wang$^{2}$, Lintao Ma$^{2}$, Zhixuan Chu$^{2}$, James Y. Zhang$^{2}$, Xiaoming Shi$^{2}$}\\\small\textbf{Pin-Yu Chen$^{3}$, Yuxuan Liang$^{6}$, Yuan-Fang Li$^{1}$, Shirui Pan$^{4}$, Qingsong Wen$^{5}$}\\[0.25em]\scriptsize $^1$Monash University, $^2$Ant Group, $^3$IBM Research, $^4$Griffith University, $^5$Alibaba Group\\\scriptsize $^6$The Hong Kong University of Science and Technology (Guangzhou)}

\makeatletter
\newcommand{\DeepPriorLogoRow}{%
  \AddToShipoutPictureFG*{%
    \AtPageUpperLeft{%
      \put(\LenToUnit{8.10in},\LenToUnit{-1.23in}){%
        \makebox[0pt][r]{
\includegraphics[width=0.825in,keepaspectratio]{deepprior-logo.png}}%
      }%
    }%
  }%
}

\newcommand{\DeepPriorTitleRow}{%
  \noindent\makebox[\textwidth][c]{%
    \begin{minipage}[c]{\textwidth}
      \centering\LARGE\bfseries \@title
    \end{minipage}%
  }\par
}

\newcommand{\DeepPriorAuthors}{%
  \begin{center}
    \begin{tabular}{c}\@author\end{tabular}
  \end{center}
}

\newcommand{\DeepPriorAbstract}[1]{%
  \begin{center}{\large\sc Abstract}\end{center}\vspace{-0.35em}%
  \begin{quote}#1\end{quote}%
}

\newcommand{\DeepPriorFigure}{%
  \noindent\makebox[\textwidth][c]{\includegraphics[width=\textwidth,keepaspectratio]{figure.png}}\par
  \vspace{2pt}%
  \noindent\makebox[\textwidth][c]{{\small Figure~1: Model Structure of Time-LLM.}}%
}

\makeatother
\begin{document}
\fancyhf{}
\renewcommand{\headrulewidth}{0pt}
\thispagestyle{empty}
\makeatletter
\DeepPriorLogoRow
\noindent\begin{minipage}[t][\dimexpr\textheight-\baselineskip\relax][s]{\textwidth}
\vspace*{0.50in}
\vfill
\DeepPriorTitleRow
\vfill
\DeepPriorAuthors
\vfill
\DeepPriorAbstract{Time series forecasting is important in many real-world dynamic systems and has been extensively studied. Unlike natural language processing and computer vision, where a single large model can tackle multiple tasks, models for time series forecasting are often specialized and require distinct designs for different tasks and applications. While pre-trained foundation models have made impressive strides in language and vision, their development in time series domains has been constrained by data sparsity. Recent studies have shown that large language models possess robust pattern recognition and reasoning abilities over complex token sequences. However, effectively aligning time series data and natural language remains challenging. This work presents Time-LLM, a reprogramming framework that repurposes large language models for general time series forecasting while keeping the backbone language model intact. Time-LLM reprograms input time series with text prototypes, uses Prompt-as-Prefix to enrich context, and projects transformed patches to obtain forecasts. Evaluations show strong performance, including few-shot and zero-shot forecasting.}
\vfill
\DeepPriorFigure
\end{minipage}
\makeatother
\end{document}